% INTRODUCTION - LLM paper profile

\subsection{Motivation}

Large language model agents are increasingly proposed as governance participants,
analysis assistants, and reporting systems in DAO operations. Yet most governance
evaluations still assume rule-based or heuristic behavior. This gap makes it difficult
to reason about where LLM augmentation improves outcomes and where it introduces new
risks, including inconsistent voting rationale, latency bottlenecks, and fragile
runtime dependencies.

\subsection{Contributions}

This paper contributes a dedicated experiment track for LLM-enabled governance:

\begin{enumerate}
    \item \textbf{Reasoning-mode comparison}: Controlled comparison of disabled, hybrid,
    and all-LLM governance modes within the same simulator.

    \item \textbf{Operational metrics}: Inclusion of LLM-specific quality signals
    (vote consistency, cache hit rate, and average latency) alongside governance outcomes.

    \item \textbf{Living-paper pipeline}: A reproducible update path that regenerates
    charts and sections as models, prompts, or agent logic change.
\end{enumerate}

\subsection{Research Questions}

This paper addresses three LLM-focused questions:

\begin{description}
    \item[RQ-L1:] How do LLM reasoning modes alter proposal pass rates, turnout, and governance activity relative to rule-based baselines?
    \item[RQ-L2:] What reliability profile emerges for LLM decision support, measured by vote consistency and cache effectiveness?
    \item[RQ-L3:] What throughput-cost tradeoffs appear when governance outcomes are analyzed against LLM inference latency?
\end{description}

\subsection{Paper Organization}

Section~\ref{sec:background} reviews related work in DAO governance and LLM agent systems.
Section~\ref{sec:theory} outlines the theoretical framing for AI-augmented governance.
Section~\ref{sec:architecture} details the simulator architecture and agent stack.
Section~\ref{sec:methodology} describes experiment design and evaluation methodology.
Section~\ref{sec:results} reports reasoning-mode outcomes and LLM systems metrics.
Section~\ref{sec:discussion} interprets implications for deployment and governance safety.
Section~\ref{sec:limitations} discusses limitations and future extensions.
Section~\ref{sec:conclusion} concludes.
